% ============================================================
% 附录 B：译注汇编（由 tools/gen_notes.py 自动生成，勿手改）
% 数据源：src/chunks/*.tex 中的 % [Tutor Note: ...] 与 \yizhu{...}
% ============================================================
\chapter{译注汇编}
\label{app:notes}

本附录收录全书的批注，共 120 条，按原书页码排列。
它们分两类。

\begin{itemize}[leftmargin=2.6em,itemsep=1pt,topsep=2pt]
  \item \textbf{校勘性批注}（原书笔误、文字层与页面图像不一致、术语取舍的理由、跨页句子的切分）。这类批注在正文中以源码注释形式存在，不出现在版面上，只在此处汇总，以免打断阅读。
  \item \textbf{概念性批注}（某个概念在通行文献里的另一种说法、与 nLab 视角的差异、值得补充的背景）。这类批注在正文中已经印出，此处再次汇总，方便通读与检索。
\end{itemize}

每条前的方括号标明它所在的章节与原书页码；页码与原书（Russian Math.
Surveys 58:3）一致，便于对照原文核查。

\section*{引言（原书 p.514）}
\addcontentsline{toc}{section}{引言（原书 p.514）}

\medskip\noindent\textbf{原书 p.514}\par
\noindent{\small\textbf{［校勘］}原书式 (1) 与 §2.1 以下各处对全导出函子的记号并不统一：本处印作 $\mathbf{R}^{\bullet}q_*$，后文（如 §2.2 定理 2.2.9 前的函子）印作 $\mathbf{R}\pi_*$。本译本一律记 $\mathbf{R}q_*$（必要时保留原书的 $\mathbf{R}^{\bullet}$），语义无别，仅记号从简。}\par\smallskip
\medskip\noindent\textbf{原书 p.515}\par
\noindent{\small\textbf{［校勘］}原书此处只写 “the simplest example of a flip or flop”，未定义 flip/flop 的一般概念；本译本保留英文原词，flip 首次出现时译作“翻转”并附原文，flop 全篇保留英文（见术语表）。}\par\smallskip

\section*{第 1 章（原书 p.517）}
\addcontentsline{toc}{section}{第 1 章（原书 p.517）}

\medskip\noindent\textbf{原书 p.517}\par
\noindent{\small\textbf{［校勘］}原文此处以括号给出同义名 exact triangle（“特异（或正合）三角”）；术语表规定全稿一律作“特异三角”，此处按原文保留括注，下文一律用“特异三角”。}\par\smallskip
\medskip\noindent\textbf{原书 p.518}\par
\noindent{\small\textbf{［概念］}公理 T1–T4 中，T1 保证恒等态射与同构不破坏特异性，T2 是旋转三角的规则，T3 使三角之间的态射可以补全，T4 通常称为八面体公理（octahedral axiom），它把两个可复合态射的锥、以及这两个锥再取的锥联系起来。本稿沿用原书编号 T1–T4，不把 T4 改称八面体公理。【来源：原书 p.517–518 的公理表述；章璞《三角范畴与导出范畴》关于八面体公理的论述（A 级）】}\par\smallskip
\medskip\noindent\textbf{原书 p.519}\par
\noindent{\small\textbf{［概念］}Serre 函子的名字来自 Serre 对偶：在光滑射影簇 $X$ 上，$\Db(\coh X)$ 的 Serre 函子由 $S(-)=(-)\otL\omega_X[\dim X]$ 给出，定义中的双函子同构正是 Serre 对偶的范畴化形式。有了它，$\Hom(A,B)$ 与反方向的 $\Hom(B,SA)$ 的对偶才能自然等同；引理 1.1.5 则说明三角范畴的 Serre 函子自动正合。【来源：原书定义 1.1.3、引理 1.1.5 及所引 [6]；章璞《三角范畴与导出范畴》中 Serre 函子与 Serre 对偶的章节（A 级）】}\par\smallskip
\medskip\noindent\textbf{原书 p.520}\par
\noindent{\small\textbf{［校勘］}页图此处把等价类的代表先写作 $(s,f)$，随即写作 $(f,s)$（原书 p.520 第 2–6 行），两处次序不一致；本译本两处均照录原文次序，未统一。}\par\smallskip
\medskip\noindent\textbf{原书 p.521}\par
\noindent{\small\textbf{［校勘］}此处原书编号印作 Definition 1.3，而非按本节编号应得的 1.1.7（本节紧接其后的命题编号为 1.1.8，可见 1.1.7 被跳过）。页图与文字层一致，均为 1.3。本译本照录原编号 1.3，未改作 1.1.7。}\par\smallskip
\noindent{\small\textbf{［概念］}$\cD/\cN$ 即 Verdier 商：对象与 $\cD$ 相同，态射由“屋顶” $X\leftarrow X'\to Y$（其锥属于 $\cN$）的等价类给出。由 $\Sigma$ 的定义，特异三角 $N\xrightarrow{\id}N\to 0\to N[1]$ 表明 $N\to 0$ 属于 $\Sigma$，故 $\cN$ 中对象在商范畴中同构于零；反之，在商范畴中同构于零的对象必属于 $\cN$。因此 $\cN$ 正是局部化函子 $Q$ 的核。【来源：原书 p.521 命题 1.1.8 前后的构造；章璞《三角范畴与导出范畴》中关于 Verdier 商的章节（A 级）】}\par\smallskip
\medskip\noindent\textbf{原书 p.522}\par
\noindent{\small\textbf{［校勘］}页图在此处把该函子印作 $\mathbf{r}FY$（$Y$ 未作下标），而在 p.521 定义处印作 $\mathbf{r}F_Y$（$Y$ 为下标）；两处同指一物。本译本统一作 $\mathbf{r}F_Y$。}\par\smallskip
\noindent{\small\textbf{［概念］}Deligne 的构造把导出函子定义为可表性：$\Rf FY$ 不由消解算出，而由反变函子 $\mathbf{r}F_Y$ 是否可表决定；等价地说，$\Rf F$ 是 $QF$ 沿局部化函子 $Q$ 的右 Kan 扩张。如此定义的导出函子对任何正合函子都有意义，既不要求 $\cC$、$\cD$ 来自阿贝尔范畴，也不要求对象有内射消解；命题 1.1.9 说明它只要在三角的两端有定义，就在整个三角上有定义。【来源：原书 p.521–522 的构造与所引 Deligne [12]（原始论文，A 级）】}\par\smallskip
\noindent{\small\textbf{［概念］}F-分裂的直观是“$F$ 在这个对象上已经算对”：典范态射 $QFY\to(\Rf F)QY$ 本身已是同构，无需再沿 $\Sigma$ 中的态射把 $Y$ 换成更好的代表。因此导出函子若要在整个范畴上有定义，每个对象都应能沿 $\Sigma$ 中的态射换成 F-分裂对象；这正是下文“足够多的 F-分裂对象”的含义，而引理 1.1.12 把 F-分裂性化归为纯可裂条件：$\Sigma$ 中任一 $s:Y\to Y'$ 的像 $QFs$ 都有收缩。【来源：原书 p.522–523 定义 1.1.11 与引理 1.1.12（A 级）】}\par\smallskip
\medskip\noindent\textbf{原书 p.523}\par
\noindent{\small\textbf{［校勘］}原书以粗体 $\mathbf{C}$、$\mathbf{H}$、$\mathbf{D}$ 区分复形范畴、同伦范畴、导出范畴，与花体 $\cC$、$\cH$ 不同；本译本照此，一律排作 $\mathbf{C}(\cA)$、$\mathbf{H}(\cA)$、$\mathbf{D}(\cA)$。}\par\smallskip
\noindent{\small\textbf{［概念］}映射锥 $Cf=M\oplus L[1]$ 的微分矩阵里，$f$ 那一项正是把 $f$ 与 $L$ 的平移耦合起来的地方。由三元组的长正合列可见，$H^{n}(Cf)$ 夹在 $H^{n}(L)$ 与 $H^{n}(M)$ 之间并与 $H^{n}(f)$ 相衔接，故 $f$ 是拟同构当且仅当 $Cf$ 零调——这正是下文定义拟同构的方式。三个箭头都取同伦类并带上横线，说明标准三角只在 $\mathbf{H}(\cA)$ 中有意义，而非 $\mathbf{C}(\cA)$ 中的序列。【来源：原书 pp.523–524；Verdier, Astérisque 239 (1996)。】}\par\smallskip
\medskip\noindent\textbf{原书 p.524}\par
\noindent{\small\textbf{［概念］}局部化 $\mathbf{H}(\cA)[\mathrm{Quis}^{-1}]$ 的含义是：把拟同构形式地变成同构，即为每个拟同构添一个逆，再对由这些逆与原态射拼出的三角关系取商。这一步之所以可行，靠的是 $\mathrm{Quis}$ 满足乘法系的 Ore 条件（1.1 节已列出局部化公理 L1–L4），于是分式可写成 $s^{-1}f$ 而不必担心分母。$\mathbf{D}^{+}$、$\mathbf{D}^{-}$、$\mathbf{D}^{b}$ 并不是这样局部化出来的商，而是 $\mathbf{D}(\cA)$ 的全子范畴，其“足够大”由引理 1.2.4 的共尾性保证。【来源：原书 pp.524–525；Gabriel–Zisman, \emph{Calculus of Fractions and Homotopy Theory}, Springer 1967。】}\par\smallskip
\noindent{\small\textbf{［概念］}这两条等价是“用内射消解计算右导出函子”的全部根据：$\mathbf{H}^{+}(\cI)$ 中的对象已是内射复形，嵌入 $\mathbf{H}^{+}(\cI)\to\mathbf{H}^{+}(\cA)$ 是全忠实的，而 $Q$ 把拟同构变成同构；两者复合为等价，等于说 $\mathbf{D}^{+}(\cA)$ 的每个对象都有内射复形代表。于是 $\mathbf{R}F$ 只须在 $\mathbf{H}^{+}(\cI)$ 上按 $F$ 直接作用即可，结果的唯一性由拟同构变成同构保证。有足够多投射对象时对偶地得到左导出函子。【来源：原书 p.524 及其所引 [17]；按所引章节编号（I.4.6、II.5.6 等）判断，[17] 当为 Hartshorne, \emph{Residues and Duality}, LNM 20, 1966。】}\par\smallskip
\medskip\noindent\textbf{原书 p.525}\par
\noindent{\small\textbf{［校勘］}引理 1.2.5 的页图与文字层一致，均印作 $\mathbf{R}FX\xrightarrow{\sim}X'$ 与 $X'\in\mathbf{H}^{+}(\cA)$。按上下文（$X\to X'$ 为拟同构，且须对 $X'$ 直接施用 $F$ 才有意义）此处当读作 $\mathbf{R}FX\xrightarrow{\sim}FX'$、$X'\in\mathbf{H}^{+}(\cR)$。本译本照录原文，仅注明于此。}\par\smallskip
\noindent{\small\textbf{［校勘］}原书 left bounded / right bounded 的用法为：left bounded 指 $p\ll 0$ 时消没（即有界在下、属 $\mathbf{D}^{+}$），right bounded 指 $p\gg 0$ 时消没（即有界在上、属 $\mathbf{D}^{-}$），与部分文献的习惯相反。本译本照原文直译为“左有界”“右有界”，读者宜以消没方向为准。}\par\smallskip
\noindent{\small\textbf{［概念］}“适应子范畴”是内射对象这一特例的公理化：条件 b) 保证每个对象都有 $\cR$-消解，条件 a) 保证这套消解在 $F$ 之下仍然零调，于是把 $F$ 作用在消解上就得到与消解选取无关的结果。内射对象关于任意加性函子都右适应，因为内射层在左正合函子下的像是零调的；所以“内射消解”只是 $\cR=\cI$ 的特例。概形上 $\cI$ 常常跑出 $\Qcoh(X)$，这时适应子范畴就取代了它。【来源：原书 p.525；同 [17]。】}\par\smallskip
\medskip\noindent\textbf{原书 p.526}\par
\noindent{\small\textbf{［校勘］}原书此处印作 “etale”，无重音符；本译本按通行拼写作 étale，并在此注明。}\par\smallskip
\noindent{\small\textbf{［概念］}三层的关系是 $\coh(X)\subset\Qcoh(X)\subset\cO_X\text{-}\Mod$。最外层是 Grothendieck 范畴，内射对象足够，但凝聚层在其中做内射消解会跑出 $\coh(X)$，所以 $\coh(X)$ 自身没有足够多内射对象；中间的 $\Qcoh(X)$ 在诺特概形上仍有足够多内射对象，又对 $\coh(X)$ 足够“大”，于是成为计算凝聚层导出函子的工作台。命题 1.3.1–1.3.3 正是在这三层之间搬运全忠实性。【来源：原书 p.526；Hartshorne, \emph{Algebraic Geometry}, GTM 52, II.5。】}\par\smallskip
\medskip\noindent\textbf{原书 p.528}\par
\noindent{\small\textbf{［校勘］}页图与文字层此处均把右导出函子的目标范畴印作 $\mathbf{D}^{+}(\cO_X\text{-}\Mod)$；由 $f:X\to Y$ 及紧随其后的“把 $\mathbf{D}^{+}(\cO_X\text{-}\Mod)_{\Qcoh}$ 送到 $\mathbf{D}^{+}(\cO_Y\text{-}\Mod)_{\Qcoh}$”一句，此处应为 $\mathbf{D}^{+}(\cO_Y\text{-}\Mod)$。本译本照录原文，仅注明于此。}\par\smallskip
\noindent{\small\textbf{［校勘］}同上一处，此式的目标范畴页图亦印作 $\mathbf{D}^{b}(\cO_X\text{-}\Mod)$，应为 $\mathbf{D}^{b}(\cO_Y\text{-}\Mod)$；照录原文。}\par\smallskip
\medskip\noindent\textbf{原书 p.529}\par
\noindent{\small\textbf{［校勘］}原书把层内的 Hom 排作花体 $\mathcal{H}om$，与整体的 $\Hom$ 相区别，其右导出函子记作 $\mathbf{R}\mathcal{H}om$、局部超 Ext 记作 $\mathcal{E}xt$；本译本保留这一花体写法，不用 preamble 中直立体的 \texttt{\textbackslash RHom}、\texttt{\textbackslash Ext}，以免混淆两者。}\par\smallskip
\noindent{\small\textbf{［概念］}$\mathcal{E}xt^{i}(\cE,\cF)$ 是 $X$ 上的层，取整体截面才得到群 $\Ext^{i}(\cE,\cF)$，两者不可混用。$\mathcal{H}om$ 只对第一个变元右正合，故 $\mathbf{R}\mathcal{H}om$ 须对第二个变元取内射消解——这解释了为什么它定义在 $\mathbf{D}(\cO_X\text{-}\Mod)^{\mathrm{op}}\times\mathbf{D}^{+}(\cO_X\text{-}\Mod)$ 上，而不是双无界。第 2 章讨论 $\mathbf{D}^{b}(X)$ 上的 $\mathcal{E}xt$ 与 Serre 对偶时，用的正是这种一边无界、一边有界的搭配。【来源：原书 p.529；同 [17]。】}\par\smallskip
\medskip\noindent\textbf{原书 p.530}\par
\noindent{\small\textbf{［校勘］}命题 1.3.6 中两个对象所属的范畴，页图与文字层一致，均印作 $\cE\in\mathbf{D}^{-}(\cO_Z\text{-}\Mod)$、$\cF\in\mathbf{D}^{-}(\cO_X\text{-}\Mod)_{\Qcoh}$。按式 (3) 的外观（$f:X\to Y$，左端 $\mathbf{R}f_{*}\cE\otL\cF$ 定在 $Y$ 上）此处应为 $\cE\in\mathbf{D}^{-}(\cO_X\text{-}\Mod)$、$\cF\in\mathbf{D}^{-}(\cO_Y\text{-}\Mod)_{\Qcoh}$；原书 $\mathcal O_Z$ 更可能是 $\mathcal O_Y$ 之误。本译本照录原文。}\par\smallskip

\section*{第 2 章（原书 p.530）}
\addcontentsline{toc}{section}{第 2 章（原书 p.530）}

\medskip\noindent\textbf{原书 p.530}\par
\noindent{\small\textbf{［概念］}Serre 函子指三角范畴 $\cD$ 上一个正合自函子 $S$，连同对任意 $A,B$ 自然的同构 $\Hom(A,B)\cong\Hom(B,S(A))^{*}$；它等价于说 $\cD$ 上存在 Serre 对偶。对光滑完备簇 $X$，$\Db(X)$ 的 Serre 函子就是 $(\,\cdot\,)\otimes\omega_X[n]$，$n=\dim X$——这正是 Serre 对偶定理在导出范畴里的整体形式，式 (6) 只是它的逐字展开。下文 (13) 处的论证（等价比与 Serre 函子交换）与本节末的 (16) 都依赖 $S$ 的存在性。〔来源：原书 §2.1 式 (6) 与定义 1.1.3；Hartshorne, \emph{Algebraic Geometry}, III.7（A 级）〕}\par\smallskip
\medskip\noindent\textbf{原书 p.531}\par
\noindent{\small\textbf{［概念］}称 $\cE$ 为函子 $\Phi_{\cE}$ 的\emph{核}（kernel object），于是 $\Phi_{\cE}$ 就是积分变换：先用 $p^{*}$ 把 $X$ 上的对象拉到积上，再与核作导出张量积，最后沿 $q$ 推回 $Y$。$\Psi_{\cE}$ 是同一核在相反方向上的变换，故 $\Phi$ 与 $\Psi$ 只在 $p$、$q$ 的次序上不同。把核取成 $\cO_{\Delta}$ 得到恒等函子；取成 Poincaré 丛则得到 Fourier–Mukai 变换——参见引言与第 5 章。〔来源：原书 §2.1 式 (7) 及其后一段；(7) 在引言式 (1) 处已出现，本处补出核的说法〕}\par\smallskip
\medskip\noindent\textbf{原书 p.532}\par
\noindent{\small\textbf{［校勘］}式 (9) 页图印作 $\mathbf{R}p_{13*}$（小写 $p$），而同一段落中三个函子印作 $\mathbf{R}\pi^{3}_{13*}$ 等（希腊 $\pi$，带上标）；文字层同样抽成 $Rp_{13*}$。这是原书记号不统一，本译本照录页图，不改成 $\pi$。}\par\smallskip
\noindent{\small\textbf{［概念］}$\cE\boxtimes_Y\cF$ 不是新的构造，按本处定义它就是沿投影拉回到三重积上再作导出张量积：$\cE\boxtimes_Y\cF:=p^{*}_{12}\cE\otL p^{*}_{23}\cF$。指标的 $Y$ 只记录“两个对象在中间的 $Y$ 方向上被粘合”这一事实。同一个符号 $\boxtimes$ 在本节后文（p.534 起）用于两个\emph{互不相干}的因子之间的外张量积，那时没有下标，含义是 $p^{*}_{13}(\cE_1)\otL p^{*}_{24}(\cE_2)$。〔来源：原书 §2.1 式 (9) 前后的定义〕}\par\smallskip
\medskip\noindent\textbf{原书 p.533}\par
\noindent{\small\textbf{［校勘］}页图此处印作 $\Aut X\times(\Pic(X)\oplus\Z)$（乘号 $\times$），而文字层抽出的是半直积号 $\ltimes$。按上下文（$G(X)$ 为半直积）应以 $\ltimes$ 为是，属原书排版笔误；本译本照录页图的 $\times$，不静默改字。}\par\smallskip
\medskip\noindent\textbf{原书 p.535}\par
\noindent{\small\textbf{［概念］}$Ad_{\cE}$ 的做法是把等价 $\Phi_{\cE}:\Db(X)\to\Db(Y)$ 与它的逆夹在两边，从而把它“共轭”成 $\Db(X\times X)$ 与 $\Db(Y\times Y)$ 之间的等价：$\cG\mapsto\Phi_{\cE}\circ\Phi_{\cG}\circ\Phi_{\cE}^{-1}$，即 (12)。它与本节前半的复合律 (9) 是同一件事的两种说法——把 $\cF\boxtimes\cE$ 当核，就得到这个共轭。$Ad_{\cE}$ 保持对角线结构层（恒等函子的核），并把 $i_{X*}\omega^{k}_{X}$ 送到 $i_{Y*}\omega^{k}_{Y}$，这正是下面造 $\mathrm{HA}(X)$ 并证明它不变的依据。〔来源：原书 §2.1 式 (11)–(13)〕}\par\smallskip
\noindent{\small\textbf{［概念］}$\mathrm{HA}(X)$ 的分次指标有两个方向：$i$ 是 $\Ext$ 的次数，$k$ 是典范层的幂次。取 $i=0$ 的一行，按典范等同它给出 $\bigoplus_{k}H^{0}(X,\omega^{k}_{X})$，即 $X$ 的典范代数（$k\geqslant0$ 的部分）与反典范代数（$k\leqslant0$ 的部分）；这就是推论 2.1.9 说 $\mathrm{HA}$ 同时含典范代数与反典范代数的原因。原书引言里这个代数印作 $\mathrm{HA}^{i,k}$（上标），§2.1 此处及 (14) 印作 $\mathrm{HA}_{i,k}$（下标），本译本各照所据页图，未作统一。〔来源：原书引言 p.515 的 $\mathrm{HA}(X)$ 公式、§2.1 本处定义与推论 2.1.9〕}\par\smallskip
\medskip\noindent\textbf{原书 p.537}\par
\noindent{\small\textbf{［校勘］}式 (18) 前后，页图与文字层都印作 $\eps\in H(X\times Y,\Q)$，缺了上标 $*$；按上下文应为 $H^{*}(X\times Y,\Q)$（(17) 处即作 $H^{*}$）。此处照录原样，不改字。}\par\smallskip
\noindent{\small\textbf{［概念］}$\vp_{\cE}$ 是函子 $\Phi_{\cE}$ 在“上同调层面”的影子：把核 $\cE$ 换成由 (18) 配出的类 $\eps$，再用 (17) 的 $v_{\eps}$ 作用。系数 $p^{*}\sqrt{\td_X}\cdot q^{*}\sqrt{\td_Y}$ 是为让复合保持复合（命题 2.1.11）而必需的修正项，其存在性由 Riemann–Roch 定理的 Grothendieck 形式保证；同一个修正项也保证推论 2.1.12 里的 K 群方块交换。$\psi_{\cE}$ 是反方向的对应物，只在需要沿 $p$ 推回时使用。〔来源：原书 §2.1 式 (18)–(19) 与命题 2.1.11 前的说明〕}\par\smallskip
\medskip\noindent\textbf{原书 p.538}\par
\noindent{\small\textbf{［校勘］}推论 2.1.12 页图与文字层均作 $\Phi_{\cE}:\Db(X)\to\Db(Z)$，而同一句里的映射印作 $\vp_{\cE}:H^{*}(X,\Q)\to H^{*}(Y,\Q)$，$Z$ 与 $Y$ 不一致；据命题 2.1.11 的记号应为 $\Db(X)\to\Db(Y)$。此处照录原样。}\par\smallskip
\noindent{\small\textbf{［校勘］}本块（ch02a）按分块约定止于 “2.2. Examples of equivalences: flopping birational transformations” 标题之前；该标题及其后内容属 ch02b。推论 2.1.12 与该 K 群交换图排在原书 p.538，故本块末页为 p.538（页数 9 页，而非分块表的 8 页），特此登记，供 Lead 校边界。}\par\smallskip
\noindent{\small\textbf{［概念］}可容许性的要点在于伴随的存在：右可容许意味着有右伴随 $P:\cD\to\cN$，于是每个对象都能“投影”到 $\cN$ 上，并由此得到特异三角 $N\to X\to M$，其余项落在正交补中。因此可容许子范畴像直和项一样把 $\cD$ 劈成两块；反过来，若对每个 $X$ 都能找到这样的三角，也可用它定义伴随。定理 2.2.8 中的 $D(X)$ 与诸 $D(Y)_k$ 都是这种意义下的可容许子范畴。【来源：原书 §2.2 定义 2.2.1–2.2.2 及其后一段；所引 [5]、[6]（A 级）】}\par\smallskip
\medskip\noindent\textbf{原书 p.539}\par
\noindent{\small\textbf{［概念］}半正交分解 $\cD=\langle\cN_0,\dots,\cN_n\rangle$ 是“带方向的直和”：条件只要求 $j<i$ 时 $\Hom(\cN_j,\cN_i)=0$，反方向的 $\Ext$ 可以非零，这正是“半”字的由来；再要求这族子范畴生成 $\cD$，每个对象便能沿方向逐层拆成特异三角。定理 2.2.8 的分解 $\langle D(Y)_{-r+1},\dots,D(Y)_{-1},D(X)\rangle$ 就是典型例子：先取 $D(X)$，再依次取扭过的 $D(Y)_k$，方向由 $k$ 的大小决定。【来源：原书 §2.2 定义 2.2.3、定理 2.2.8；所引 [5]、[6]（A 级）】}\par\smallskip
\noindent{\small\textbf{［概念］}例外对象 $E$ 的自同态代数只有 $\Hom^0(E,E)=k$，高次 $\Ext$ 全消没，故它在导出范畴里像一点一样刚性；一族例外对象若两两之间只有单向的 $\Ext$（$i>j$ 时 $\Hom^{\bullet}(E_i,E_j)=0$）且生成整个范畴，就给出该范畴与 $\End(\bigoplus_i E_i)$ 上有限维模的导出范畴的等价。$\PP^N$ 上的 $(\cO(i),\dots,\cO(i+N))$ 是最早也最典型的例子。【来源：原书 §2.2 定义 2.2.4、例 2.2.5 及所引 [2]（A 级）】}\par\smallskip
\medskip\noindent\textbf{原书 p.540}\par
\noindent{\small\textbf{［校勘］}页图此处印作 “the projective bundle of the exceptional divisor of $\tilde Y$ over the centre $Y$”，字面上是“例外除子 $\tilde Y$ 的射影丛”；按上下文，$\tilde Y$ 本身就是爆破 $\tilde X\to X$ 的例外除子，故译作“例外除子 $\tilde Y$ 在爆破中心 $Y$ 上的射影丛”。}\par\smallskip
\noindent{\small\textbf{［校勘］}页图此式末项印作 $\Rf\pi_*\cO_{\tilde Y}\otL G$，而同一证明下一行印作 $\Rf\pi_*\cO_{\tilde X}\cong\cO_X$；按投影公式应为 $\cO_{\tilde X}$，原书此处当作 $\tilde X$。本译本照录页图，读者请按 $\tilde X$ 理解。}\par\smallskip
\medskip\noindent\textbf{原书 p.541}\par
\noindent{\small\textbf{［校勘］}p.541 图 (20) 的页图由缩微胶片扫描而来，底排两个节点（$X$ 与 $X^+$）的标签、以及右下角若干残段均已糊死，无法逐字辨认；页图右端残留的短虚线亦不能判定归属。本图按正文所述映射重绘：6 个节点为 $\tilde Y$、$Y$、$\tilde X$、$Y^+$、$X$、$X^+$，7 条实线箭头为 $p:\tilde Y\to Y$、$j:\tilde Y\to\tilde X$、$p^+:\tilde Y\to Y^+$、$i:Y\to X$、$\pi:\tilde X\to X$、$\pi^+:\tilde X\to X^+$、$i^+:Y^+\to X^+$。凡页图不能证实的虚线箭头（如以 $\mathrm{fl}$ 标注的 $X\dashrightarrow X^+$）一概不画，以免杜撰。}\par\smallskip
\noindent{\small\textbf{［概念］}定理 2.2.8 就是通常所说的 Orlov 爆破定理：$\Db(\tilde X)$ 由 $\Db(X)$ 与若干份“扭过的” $\Db(Y)$ 拼成，份数为余维 $r$ 减一。写成半正交分解后，$\Db(\tilde X)$ 的每个对象都能沿 $D(X),D(Y)_{-r+1},\dots$ 逐层拆开，于是关于爆破的问题可化归为关于 $X$ 与 $Y$ 的问题；本节后文正是用它把对象 $\overline A$ 定位到 $D(X)^{\perp}=\langle D(Y)_{-l},\dots,D(Y)_{-1}\rangle$ 之中。【来源：原书 §2.2 定理 2.2.8 及所引 [34]、[7]（A 级）】}\par\smallskip
\noindent{\small\textbf{［概念］}flip 与 flop（翻模）都是双有理变换，都把 $X$ 中的子簇 $Y$ 换成 $X^+$ 中的另一子簇 $Y^+$，且被改动的子簇余维 $\geqslant2$。区别在典范类：flip 两侧的典范类不同，收缩曲线上的 $K$ 由负变正；flop 两侧的典范类相同，$K_{X}$ 与 $K_{X^+}$ 在公共的爆破 $\tilde X$ 上重合。本例正体现这一点：$l<k$ 时 $Y=\PP^k$ 被换成维数更低的 $Y^+=\PP^l$，是 flip；$l=k$ 时 $Y$ 与 $Y^+$ 都是 $\PP^k$，而式 (21) 中 $\omega_{\tilde X}|_{\tilde Y}=\cO(-k;-k)$ 关于 $\tilde Y=\PP^k\times\PP^k$ 的两个因子对称，是 flop。【来源：原书 §2.2（$l<k$ 与 $l=k$ 的区分）与式 (21)；flip、flop 的一般定义见 Kollár–Mori《Birational Geometry of Algebraic Varieties》（A 级）】}\par\smallskip
\medskip\noindent\textbf{原书 p.543}\par
\noindent{\small\textbf{［校勘］}文字层把 (23)、(24) 及以下各处的 $\overline A$ 一律抽成 $A$（上横线丢失）；页图三处均印作 $\overline A$。本译本从页图，全篇作 $\overline A$，以免与证明开头的对象 $A\in\Db(X^+)$ 混淆。}\par\smallskip
\medskip\noindent\textbf{原书 p.544}\par
\noindent{\small\textbf{［校勘］}本块止于此处：其下即为第 3 章标题 “CHAPTER 3 — Fully faithful functors between derived categories”，属 ch03a 的内容。}\par\smallskip

\section*{第 3 章（原书 p.544）}
\addcontentsline{toc}{section}{第 3 章（原书 p.544）}

\medskip\noindent\textbf{原书 p.544}\par
\noindent{\small\textbf{［校勘］}原书节标题与正文接排（“3.1. Postnikov diagrams and their convolutions. In this section…”），本译本按全稿体例改为独立节标题，并加节首重点框。}\par\smallskip
\noindent{\small\textbf{［校勘］}页图上斜箭头的落点在 200 dpi 下压着底行字缘，无法直读；此处按 $j_{k}\circ i_{k}=d^{k}$ 与 $Y^{c}=X^{c}$、$i_{c}=\id$ 的相容性判定为 $i_{k}:X^{k}\to Y^{k}$、$j_{k}:Y^{k}\to X^{k+1}$，与原文同段所述的交换三角一致。}\par\smallskip
\medskip\noindent\textbf{原书 p.545}\par
\noindent{\small\textbf{［概念］}Postnikov 系统是三角范畴中对“全复形”的替代：三角范畴没有核与余核，无法像在复形范畴里那样直接取 $\Tot$，这里的做法是逐级把相邻两项的微分换成映射锥，得到一串首尾相接的特异三角，末项 $Y^{0}$ 就是卷积。同一个复形可以有互不同构的卷积，$\Tot(X^{\cdot})$ 因此一般只是一个类，既可能含多个对象，也可能为空。〔译者按（上下文推断，无外部依据）：图中圆圈标出的是 $j_{k}\circ i_{k}=d^{k}$ 这一交换性，正是它保证逐级的锥与原微分相容。〕【来源：原书 §3.1，p.544–545】}\par\smallskip
\medskip\noindent\textbf{原书 p.546}\par
\noindent{\small\textbf{［校勘］}文字层把 $\overline{d^{c+1}}$ 的上划线整个丢掉，抽成 $d^{c+1}$（例如“there is a unique morphism dc+1 : Y c+1 → Xc+2 such that dc+1 ◦ $\alpha$ = dc+1”），页图印作 $\overline{d^{c+1}}$；本译本按页图加上划线，以便与同段中的 $d^{c+1}$ 区分。}\par\smallskip
\medskip\noindent\textbf{原书 p.547}\par
\noindent{\small\textbf{［校勘］}式 (30) 页图印作 $\mathbf{R}\pi_{*}$（$\mathbf{R}$ 上没有点），而引言式 (1) 同一类函子印作 $\mathbf{R}^{\bullet}q_{*}$；本译本照录式 (30) 的写法。}\par\smallskip
\noindent{\small\textbf{［概念］}形如 (30) 的函子即所谓\emph{积分变换}：核对象 $\cE$ 取自积 $M\times X$，函子先沿 $p$ 拉回、再与 $\cE$ 作导出张量积、最后沿 $\pi$ 前推。上一段末尾的两个公式给出它的左、右伴随，其中 $\cE^{\vee}$ 是 $\cE$ 的对偶复形。定理 3.2.1 断言：只要正合函子全忠实且有伴随，它就一定是某个核对象的积分变换，且核在同构意义下唯一——这正是第~3 章要证的“等价总可由积上的对象表示”。【来源：原书引言式 (1) 与 §3.2 式 (30)、p.547–548】}\par\smallskip
\medskip\noindent\textbf{原书 p.548}\par
\noindent{\small\textbf{［概念］}Serre 函子把 Serre 对偶范畴化：在光滑射影簇 $X$ 的凝聚层有界导出范畴上，它由 $S_{X}=(\,\cdot\,)\otimes\omega_{X}[\dim X]$ 给出（原书所引式 (6) 即此）。有了它，左伴随 $F^{*}$ 一旦存在，右伴随便由 $F^{!}=S_{M}\circ F^{*}\circ S_{X}^{-1}$ 给出，无须另行构造——这正是“有左伴随”与“有右伴随”可互相推出的依据。【来源：原书 p.548 及所引式 (6)】}\par\smallskip
\noindent{\small\textbf{［概念］}有界性说的是：$F$ 把 $\Db(M)$ 中对象的上同调整体压进一个长度固定为 $n$ 的区间 $[z,z+n]$，两个端点都不随对象改变。引理 3.2.3 表明，只要 $F$ 有左伴随，这一性质自动成立：证明取极丰富线丛 $\cL$，用 $\cO(k)$ 的右消解与一个谱序列控制住 $G(\cL^{k})$（$k<0$）的上同调，再由 $\cL$ 的丰富性回到 $F$ 本身。【来源：原书 §3.2，p.548–549】}\par\smallskip
\medskip\noindent\textbf{原书 p.549}\par
\noindent{\small\textbf{［校勘］}谱序列右端页图与文字层都印作 $G(\cL^{i})$，而同一句的结论写作 $G(\cL^{k})$；据上下文此处应为 $k$，但照录原文，不改。}\par\smallskip
\noindent{\small\textbf{［校勘］}本块止于注 3.2.4；“3.3. Construction of the object representing a fully faithful functor”标题及其后内容归 ch03b，此处不译。}\par\smallskip
\noindent{\small\textbf{［概念］}式 (31) 就是 Beilinson 为 $\PP^{N}\times\PP^{N}$ 的对角所造的 Koszul 型消解：第 $-i$ 项取 $\cO(-i)\boxtimes\Omega^{i}(i)$，微分由外积与缩并的配对给出，整条复形只用到 $\PP^{N}$ 上的欧拉序列。把 $\Omega^{i}(i)$ 逐项换成 $\bigwedge^{i}A_{1}\otimes\cO$，就得到后文 (35) 中射影空间上的消解；两者是同一个 Koszul 复形的几何版与代数版。【来源：原书 p.549 式 (31) 及 p.553 式 (35)。】}\par\smallskip
\medskip\noindent\textbf{原书 p.550}\par
\noindent{\small\textbf{［概念］}这里所谓复形 $C^{\bullet}$ 的卷积，是把各项按次数平移后拼成一个对象：取 $\bigoplus_{i}C^{i}[-i]$，微分取 $d_{C}$ 与平移带来的符号 $(-1)^{i}$ 之和；其存在即引理 3.1.2 的结论，该引理同时保证，当各项之间的负次数 $\Hom$ 全为零时这样的对象在同构意义下唯一。后文反复使用这一点：只要认出两个对象是同一个复形的卷积，就得到它们同构。【来源：原书 p.550 对引理 3.1.2 的援引，及 p.551、p.553 的同类论证。】}\par\smallskip
\medskip\noindent\textbf{原书 p.551}\par
\noindent{\small\textbf{［校勘］}页图此句印作 $\pi'_{*}(C^{\bullet}_{k})$，没有粗体 $\mathbf{R}$；而同一句后半与 p.554 以下一律印 $\mathbf{R}\pi'_{*}$。此处照录页图，读者按导出的直接像理解即可。}\par\smallskip
\medskip\noindent\textbf{原书 p.552}\par
\noindent{\small\textbf{［校勘］}页图此处印作 “the existence of an object in the category $E\in\mathbf{D}^{b}(M\times X)$”，把对象名 $E$ 与范畴名错位；页图与文字层一致，非抽字错误。译文按原意作“范畴 $\Db(M\times X)$ 中存在对象 $\cE$”。}\par\smallskip
\noindent{\small\textbf{［概念］}$n$-Koszul 性是 Koszul 性的有限截断：只要求前 $n$ 步的模序列正合，不问整条消解。由 (33) 递归给出的 $B_{m}$ 正是 Koszul 复形的各项；p.553 的自然嵌入 $\bigwedge^{i}A_{1}\subset B_{i}$ 说明 $B_{i}$ 至少含外代数的线性部分，这正是 Koszul 性在交换情形下的含义。【来源：原书 p.552 定义 3.3.3、p.553 的嵌入，以及文献 [1]（Backelin）关于 Veronese 代数的结果。】}\par\smallskip
\medskip\noindent\textbf{原书 p.554}\par
\noindent{\small\textbf{［校勘］}页图此图下行的左端印作 $\mathbf{R}j_{*}\cL^{-n}$ 与 $\mathbf{R}j_{*}\cO_{M}$（小写 $j$），而 p.549、p.552、p.555 处同一个函子一律印 $\mathbf{R}J_{*}$，其中 $J=(j\times\id)$。据上下文此处应为 $\mathbf{R}J_{*}$；此处照录页图，不改字。}\par\smallskip
\medskip\noindent\textbf{原书 p.555}\par
\noindent{\small\textbf{［校勘］}页图此图左上印 $S^{k}A_{1}\otimes F(\cO)$，左下印 $A_{k}\otimes F(\cO)$，二者并置。这并不是笔误：当分次代数 $A$ 由 $A_{1}$ 生成时有 $S^{k}A_{1}=A_{k}$，此处照录页图。}\par\smallskip
\noindent{\small\textbf{［校勘］}页图此条末尾 “unique up to isomorphism” 之后无句点，与全书陈述末尾不加标点的排法一致；译文按中文习惯仍以句号收尾。}\par\smallskip
\noindent{\small\textbf{［概念］}整节的路线可以概括为“先放大、再切掉”：借极丰富线丛把 $M$ 嵌进 $\PP^{N}$，使 $M$ 上的层都能用 $\cO(k)$ 的消解表示，于是 $F$ 的作用就转写成 $\PP^{N}\times X$ 上的对象 $\cE'$；再用 (36) 的卷积 $G$ 与同调维数估计把多余的部分 $C$ 分离出去，剩下 $\cE$。$K$ 与 $\cE'$ 的差别由特异三角 $S\to K\to\cE'\to S[1]$ 度量，而 $S$ 的上同调集中在 $[-n-a,-n]$，正好不与 $\Rf J_{*}\cE$ 所在的区间 $[-a,0]$ 相交。【来源：原书 p.549–555 全节结构及 p.555 末的论证。】}\par\smallskip
\noindent{\small\textbf{［概念］}丰富序列（ample sequence）是这里的技术核心：它把“导出范畴中的对象为零”这一整体性断言，化归为“在一列可数对象上所有阶的 $\Hom$ 全为零”这一逐项可检验的条件（引理 3.4.3、3.4.4）。条件 a) 与 b) 保证序列中的对象能“生成”任意对象 $X$ 并给出消解，条件 c) 则保证 $\Hom(X,P_i)$ 一侧的消没；三者合起来使 $P_i$ 在指标充分负时表现得像 $\PP^n$ 上的 $\cO(-N)$ 那样“足够负”。【来源：原书 §3.4 定义 3.4.1 及紧随其后的例 3.4.2（A 级）】}\par\smallskip
\medskip\noindent\textbf{原书 p.556}\par
\noindent{\small\textbf{［校勘］}页图把分次 Hom 的上标印作小圆点（Hom$^{\cdot}$），文字层抽作“Hom·”；本译本统一记作 $\Hom^{\bullet}$，语义不变。}\par\smallskip
\noindent{\small\textbf{［概念］}这里 $\Hom^{\bullet}$ 指分次 $\Hom$，即 $\Hom^{\bullet}(E,F)=\bigoplus_{k\in\Z}\Hom^{k}(E,F)$；等式 $\Hom^{\bullet}(P_i,X)=0$ 的意思是所有阶的 $\Hom$ 同时为零，而非某一阶为零。页图把上标印成一个小圆点（$\Hom^{\cdot}$），本译本统一记作 $\Hom^{\bullet}$。【来源：原书 p.556 引理 3.4.3 及其后各处的用法（A 级）】}\par\smallskip
\medskip\noindent\textbf{原书 p.558}\par
\noindent{\small\textbf{［校勘］}页图此处中间一条竖直箭头的标签印作 $v_2$；按上下文，它是同构 $f_{i_1}:P_{i_1}^{\oplus k_1}\iso F(P_{i_1}^{\oplus k_1})$，与 $f_j$、$f_X$ 并列，页图误排为 $v_2$。此处照录。}\par\smallskip
\medskip\noindent\textbf{原书 p.561}\par
\noindent{\small\textbf{［校勘］}页图此式第四项印作 $F(\varphi)F_{Y_j}y$；按上下文，此处应是 $F(\varphi)f_{Y_j}y$（$f_{Y_j}$ 为归纳假设给出的同构），页图漏掉了下标小写 $f$ 而误排为 $F_{Y_j}$。此处照录。}\par\smallskip
\medskip\noindent\textbf{原书 p.562}\par
\noindent{\small\textbf{［校勘］}页图 (41) 末项印作 $P_i^{\oplus k}$；按特异三角的写法，第三支的终点应是 $P_i^{\oplus k}[1]$。此处照录。}\par\smallskip
\noindent{\small\textbf{［校勘］}页图 (42) 中间一项印作 $f(\beta\varphi)f_X$（小写正体 $f$ 后直接接括号），据上下文应为 $F(\beta\varphi)f_X$。此处照录。}\par\smallskip
\medskip\noindent\textbf{原书 p.563}\par
\noindent{\small\textbf{［校勘］}页图此处印作“应用引理 3.6”；本章并无编号为 3.6 的引理，所指应是引理 3.4.3（丰富序列上 $\Hom^{\bullet}$ 全为零则对象为零）。原文照录，未改编号。}\par\smallskip
\noindent{\small\textbf{［校勘］}页图此段印作 $F^{!}(Z)=0$；按上下文，$F^{!}(f_Z)$ 是同构只能推出映射锥满足 $F^{!}(C_Z)=0$（随即用于 $\Hom(F(Y),C_Z)=0$），故原式应为 $F^{!}(C_Z)=0$。此处照录。}\par\smallskip
\noindent{\small\textbf{［概念］}这里用了两次“三角分裂”：映射锥给出特异三角 $F(Z)\longrightarrow\Phi_{\cE}(Z)\longrightarrow C_Z\longrightarrow F(Z)[1]$。由 $\Hom(\Phi_{\cE}(Z),C_Z)=0$，该三角分裂，于是 $F(Z)\cong C_Z[-1]\oplus\Phi_{\cE}(Z)$，平移得 $F(Z[1])\cong C_Z\oplus\Phi_{\cE}(Z)[1]$；原文随即把 $Y=Z[1]$ 代入已证的 $\Hom(F(Y),C_Z)=0$，而该直和项含 $\id_{C_Z}$，故 $C_Z=0$。〔译者按（上下文推断，无外部依据）：分裂后的直和分解与代入 $Y=Z[1]$ 这一步，原书未写出，为译者补出。〕}\par\smallskip
\noindent{\small\textbf{［概念］}“复形 (32) 的卷积”指把该复形各阶项按 Postnikov 系统逐步拼接成单个对象；同一复形的不同卷积之间一般并不同构，此处的关键恰是引理 3.1.2 断言本复形的所有卷积彼此同构（这正是 §3.1 中 Postnikov 图技术的用途）。【来源：原书 §3.1 引理 3.1.1–3.1.2 与 §3.3 命题 3.3.5（A 级）】}\par\smallskip
\medskip\noindent\textbf{原书 p.564}\par
\noindent{\small\textbf{［校勘］}文字层把花体的 $\mathcal{J}$ 一律抽成 $I$（如“I(m) k”）；页图印作 $\mathcal{J}^{(m)}_k$，本译本从页图作 $\cJ$。}\par\smallskip
\noindent{\small\textbf{［校勘］}页图与文字层此处均印作嵌入到 $X^n$；但 $s$ 应是 $S^{(m)}_k\subset X^m$ 到 $X^m$ 的嵌入（引理 3.5.1 的各项都在 $X^m$ 上取上同调），据上下文应为 $X^m$。此处照录。}\par\smallskip
\noindent{\small\textbf{［概念］}齐次坐标代数（homogeneous coordinate algebra）指射影簇 $X$ 关于极丰富线丛 $\cL$ 的齐次坐标环：把 $\cL$ 各次幂的截面按次数汇成分次代数 $A=\bigoplus_{k\geqslant 0}H^0(X,\cL^k)$。当 $\cL=\cO_X(1)$ 由 $X\hookrightarrow\PP^N$ 给出、且 $X$ 射影正规时，$A$ 就是 $X$ 在 $\PP^N$ 中的齐次坐标环。分次结构使“以 $\cL$ 为权”的上同调能按次数逐块读取，这正是下文把 $A_m=H^0(X,\cL^m)$ 当作系数使用的原因；附加条件 $H^i(X,\cL^k)=0$（$k>0$，$i\ne 0$）则排除了正次数上的高阶上同调。【来源：原书 §3.5 开头及 §3.3 式 (33) 前；Polishchuk–Positselski, \emph{Quadratic Algebras}, AMS Univ. Lecture Ser. 37, 2005, 第 1 章（A 级，经典教材）】}\par\smallskip
\medskip\noindent\textbf{原书 p.565}\par
\noindent{\small\textbf{［校勘］}本命题即 §3.3 命题 3.3.4 的“证明版本”：把那里的簇 $M$ 改写为 $X$，并显式给出条件 (43)。第 3) 条中的 $\cO_M$ 沿用了 §3.3 的记号——那里规定 $R_0=\cO_M$；本节簇改称 $X$，故 $\cO_M$ 与 $\cO_X$ 是同一个层，并非笔误。此处照录。}\par\smallskip
\noindent{\small\textbf{［概念］}Koszul 复形（Koszul complex）原指由环中元素序列 $f_1,\dots,f_n$ 以 $\bigwedge^\bullet$ 为项作出的复形：它的正合性等价于该序列构成正则序列，这是“线性消解”的原型。本节的 $P^\bullet_n$ 取同样的形状，项换成理想层 $J_{\Sigma^{(n)}}$ 与 $\cJ^{(n)}_i$，整个复形在 $X^n$ 上正合，给出 $J_{\Sigma^{(n)}}$ 的一个消解；$n=2$ 时它退化为 $0\to J_\Delta\to\cO_{X\times X}\to\cO_\Delta\to 0$。附录正是借它把“对角有消解”与“$A$ 是 n-Koszul 代数”两件事接起来。【来源：原书 §3.5；Koszul 复形的经典表述见 Weibel, \emph{An Introduction to Homological Algebra}, CUP 1994, §4.5（A 级）】}\par\smallskip
\noindent{\small\textbf{［概念］}n-Koszul 性质的条件已由 §3.3 定义 3.3.3 给出：$k$ 作为右 $A$-模有自由消解 $B_n\otimes_k A\to\cdots\to B_1\otimes_k A\to A\to k\to 0$，其中 $B_1=A_1$，$B_m=\Ker(B_{m-1}\otimes A_1\to B_{m-2}\otimes A_2)$ 逐次取核而得（式 (33)）。$n=2$ 时这就是经典的 Koszul 代数，一般 $n$ 只是把“线性消解”截到第 $n$ 步。判断 Koszul 性的标准不变量是 Ext 代数 $\Ext^\bullet_A(k,k)=\bigoplus_i\Ext^i_A(k,k)$：$A$ 为 Koszul 代数，当且仅当这一分次代数由 $\Ext^1_A(k,k)$ 生成。附录不去计算 Ext 代数，而是直接给出使上述消解正合的充分条件 (43)。【来源：原书 §3.3 定义 3.3.3 与式 (33)、§3.5 命题 3.5.3；Polishchuk–Positselski, \emph{Quadratic Algebras}, AMS Univ. Lecture Ser. 37, 2005, 第 1–2 章（A 级，经典教材）】}\par\smallskip
\noindent{\small\textbf{［概念］}对角的消解（resolution of the diagonal）指 $X\times X$ 上以结构层 $\cO_\Delta$ 为末项的复形。它的用处在于：把 $\cO_\Delta$ 换成由外张量积 $\cL^a\boxtimes\cL^b$ 这类简单对象组成的复形之后，关于“函子的核”的论断就可以逐项验证。命题 3.5.3 第 3) 条给出的正是一条长度 $n$ 的这类消解 $\cL^{-n}\boxtimes R_n\to\cdots\to\cO_M\boxtimes R_0\to\cO_\Delta$（§3.3 规定 $R_0=\cO_M$），它与 $A$ 的 n-Koszul 消解互为表里；§3.3 在 $\PP^N\times\PP^N$ 上先用的是 Beilinson 型消解（文献 [2]），附录补上 $X\times X$ 上与之配套的 n-消解。【来源：原书 §3.3 命题 3.3.4 及其后的用法、§3.5 命题 3.5.3 与式 (46)】}\par\smallskip
\medskip\noindent\textbf{原书 p.566}\par
\noindent{\small\textbf{［校勘］}页图印作 “the exactness for $i\leqslant 1$”，而上一句的序列是对 $i\geqslant 1$ 给出的；按上下文这里指 $i=1$ 的情形。此处照录。}\par\smallskip
\medskip\noindent\textbf{原书 p.567}\par
\noindent{\small\textbf{［校勘］}式 (46) 与命题 3.5.3 第 3) 条一样沿用 $\cO_M$（即本节语境下的 $\cO_X$），见前注；此处照录。}\par\smallskip
\noindent{\small\textbf{［校勘］}本块自 p.555 中部起。p.555 上半页为命题 3.3.5 的证明结尾与其陈述，归 ch03b；本块仍标 p.555 页标，以便页码覆盖连续。}\par\smallskip

\section*{第 4 章（原书 p.568）}
\addcontentsline{toc}{section}{第 4 章（原书 p.568）}

\medskip\noindent\textbf{原书 p.568}\par
\noindent{\small\textbf{［校勘］}文字层把 Mukai 格记号的波浪线整体丢失（抽出为 “H(S, Z)”），页图作 $\tilde H(S,\Z)$；本译按页图补出。又原书该记号用直立罗马体 H，本稿为与全书的 $H^*(S,\Z)$ 一致改用斜体 $H$ 加波浪线，差异登记在报告 §5。}\par\smallskip
\noindent{\small\textbf{［概念］}这组数值与上一句的格论描述互为表里：$H^2(S,\Z)$ 上相交型的符号差是 $(3,19)$，故秩为 22；Noether 公式给出 $\chi(\cO_S)=2$，与 $c_2(S)=24$ 一起决定 $p_g=1$、$h^{1,1}=20$。原书只用一句带过，读者若初次接触 K3 曲面，宜把这三个数当作同一批事实的不同写法。【来源：原书 p.568；Barth–Hulek–Peters–Van de Ven《Compact Complex Surfaces》中 K3 曲面一章（A 级）】}\par\smallskip
\noindent{\small\textbf{［概念］}配对中 $-\alpha\cdot\alpha'$ 的负号来自 Serre 对偶：$\Ext^2$ 的维数由 $\Hom$ 的对偶给出，于是 $\chi$ 在中间次数上带负号。$\mathrm{U}$ 因子对应 $H^0\oplus H^4$ 上的双曲配对，正是它使两类初等等距（线丛陈特征的平移与 $r$、$s$ 的互换，见 p.571）成为 Mukai 格的自等距，后文归约 $r>1$ 全靠这一条。【来源：原书 p.568 与 p.571 的 $\varphi_m$ 公式（A 级）】}\par\smallskip
\medskip\noindent\textbf{原书 p.569}\par
\noindent{\small\textbf{［校勘］}命题 4.1.3 第 1) 条把所构造的等距记作 $f_{Z_{\cE}}$，第 2) 条起改记作 $f$，本页证明与 p.570 也一律用 $f$；原书记号不一致，照录。}\par\smallskip
\medskip\noindent\textbf{原书 p.570}\par
\noindent{\small\textbf{［校勘］}上面三行等式的第二个等号把 $\nu_{S_2,*}\pi_*$ 直接换成了 $\nu_{S_1\times S_2,*}$，理由是 $\nu_{S_2}\circ\pi=\nu_{S_1\times S_2}$；原文未写出这一步，属跳步。}\par\smallskip
\medskip\noindent\textbf{原书 p.571}\par
\noindent{\small\textbf{［校勘］}此处“由于闭链 $Z$ 是代数的”中的 $Z$ 应指 $Z_{\cE}$——本节先前只定义了 $Z_{\cE}$，未定义过 $Z$；照录原文记号。}\par\smallskip
\noindent{\small\textbf{［校勘］}$\varphi_m$ 中的 $(m,l)$ 是 $H^2$ 上通常的相交配对，而 p.568 规定 Mukai 格在 $H^2$ 上取负号；原书此式按相交型写出，与前文的符号约定并存，读者套用公式时须留意。}\par\smallskip
\noindent{\small\textbf{［概念］}这一段的三项条件各司其职，缺一不可：$v$ 迷向保证 $\Ext^1$ 消没，从而不同点的层彼此正交；$v$ 落在正锥内保证模空间非空（这一点原文由 Mukai 的定理一并引用，未展开）；$(v,u)=1$ 给出精细模空间，积上才有泛对象 $\cE$。若只有迷向而没有最后这个本原性条件，得到的只是粗模空间，$\Phi_{\cE}$ 便无从定义。【来源：原书 p.571 及其所引 Mukai [31]（A 级）】}\par\smallskip
\medskip\noindent\textbf{原书 p.572}\par
\noindent{\small\textbf{［校勘］}前一句写 $f^{-1}\circ g$，紧接着写 $f^{-1}\cdot g$，同一个复合在两处记号不同；照录原文。}\par\smallskip

\section*{第 5 章（原书 p.573）}
\addcontentsline{toc}{section}{第 5 章（原书 p.573）}

\medskip\noindent\textbf{原书 p.573}\par
\noindent{\small\textbf{［校勘］}定义 5.1 前句用 $A_1,\dots,A_m$ 与 $(\alpha_1,\dots,\alpha_m)$，后句却写 $\cP_{(\alpha_1,\dots,\alpha_k)}$ 与 $A_1\times\cdots\times A_k$；两处下标应统一为 $m$。原书如此，照录。}\par\smallskip
\noindent{\small\textbf{［校勘］}本段页图把 $\hat A\times\widehat{\hat A}$ 上的 Poincaré 丛印作斜体 $P_{\hat A}$，而同一句的 $\cP_A$ 用花体 $\cP$；同一丛出现两种字体，照录。}\par\smallskip
\noindent{\small\textbf{［概念］}这里的“细”是关键：$\hat A$ 是细模空间，故存在泛线丛；再要求它在 $\{e\}\times\hat A$ 上平凡，才把 $\cP$ 唯一钉住——否则还可以乘上任一从 $\hat A$ 提升的线丛。命题 5.1.2 说相应的 $\Phi_\cP$ 是等价，这正是 Mukai 的原始结果，也是本节全部构造的起点。【来源：原书 p.573–574 与所引 [29]（Mukai，A 级）】}\par\smallskip
\noindent{\small\textbf{［概念］}$A$ 与 $\hat A$ 一般并不同构，它们的导出范畴却等价：这是历史上第一个“不同簇、同范畴”的例子（见引言）。$\Phi_\cP$ 就是 Fourier–Mukai 变换的原型，$A\times\hat A$ 上的 Poincaré 丛是它的核；命题 5.1.2 中的 $\Psi_\cP$ 则是它的逆（差一个平移与群逆）。【来源：原书引言（p.513）与 §5.1（p.573–574），并见所引 [29]（Mukai，A 级）】}\par\smallskip
\medskip\noindent\textbf{原书 p.574}\par
\noindent{\small\textbf{［校勘］}页图此处印 $p^*_1(\cP_\alpha)$，而 (50) 与命题 5.1.5 b) 均印 $p^*_2$；按 (50) 应为 $p^*_2$。原书如此，照录。}\par\smallskip
\medskip\noindent\textbf{原书 p.575}\par
\noindent{\small\textbf{［校勘］}引理 5.1.8 末字页图印作带上划线的 $\subset$（即 $j=\overline{\subset}$），当是 $j=\bar\iota$ 之误排，其中 $\iota$ 为 $i:Y\into X$ 在 $\bar k$ 上的基变换；语义为 $j$ 是 $i$ 的基变换。原书如此，照录。}\par\smallskip
\noindent{\small\textbf{［概念］}两个引理合起来是一次下降论证：引理 5.1.8 把“闭子簇上的可逆层”沿域扩张 $\bar k/k$ 下降回 $k$，引理 5.1.9 说明全忠实与等价在域扩张下不变。于是定理 5.1.7 的证明可设 $k$ 代数闭，第 1–3 步的几何论证才得以展开；否则摩天楼层 $\cO_e$、仿射邻域与极大理想 $\mathfrak m$ 这些工具都要换成一般域上的说法。【来源：原书 p.575–576 与所引 [37]（A 级）】}\par\smallskip
\medskip\noindent\textbf{原书 p.577}\par
\noindent{\small\textbf{［校勘］}第 3 步两处把 $(a,\alpha)$ 印作 $(a,a')$：一处在“写成和 $(a,a')=(a_1,\alpha_1)+(a_2,\alpha_2)$”，一处在下段“对任意闭点 $(a,a')$”。按上下文应为 $(a,\alpha)$。原书如此，照录。}\par\smallskip
\noindent{\small\textbf{［校勘］}上条同构链中第二、三个 $\cong$ 各只保留复合的一个因子；按 $\Phi_{(a,\alpha)}=\Phi_{(a_1,\alpha_1)}\Phi_{(a_2,\alpha_2)}$，前两步应合写为 $\Phi_\cE\circ\Phi_{(a_1,\alpha_1)}\Phi_{(a_2,\alpha_2)}\circ\Phi_\cE^{-1}$。原书如此，照录。}\par\smallskip
\noindent{\small\textbf{［校勘］}原书说 $\Phi_{(a,\alpha)}$ “按定义 (49) 等于 $T^*_a(\,\cdot\,)\otimes\cP_\alpha$”，但 (49) 印的是 $T^*_{-a}$；且 $T^*_a=T_{-a*}$，与 (49) 的 $T_{a*}$ 并不同。此处应为 $T^*_{-a}$。原书如此，照录。}\par\smallskip
\medskip\noindent\textbf{原书 p.578}\par
\noindent{\small\textbf{［校勘］}命题 5.1.11 陈述中前一处印 $\cP_{(\alpha,-a)}$，后一处却印 $\cP_{(\alpha,a)}$；又“5.1.1”应是“定义 5.1”。两处下标与编号均疑为原书笔误，照录。}\par\smallskip
\medskip\noindent\textbf{原书 p.580}\par
\noindent{\small\textbf{［概念］}“等距”一词来自 $A\times\hat A$ 上由 Poincaré 丛 $\cP$ 给出的结构：条件 $\tilde f=f^{-1}$ 把 $f$ 的矩阵与它的对偶矩阵（转置并变号）绑在一起，正是“保持这一结构”的写法。命题 5.1.15 表明，任何导出等价给出的 $f_\cE$ 都自动等距，因此定理 5.1.16 所得的同构比一般的簇同构强。【来源：原书 p.579–580（定义 5.1.14、命题 5.1.15）与所引 [38]（A 级）】}\par\smallskip
\medskip\noindent\textbf{原书 p.581}\par
\noindent{\small\textbf{［校勘］}文字层把 $\mathbf{R}^{\bullet}\mathcal{H}om$ 抽成 “R·Hom”，页图印作 $\Rb\cH\mathrm{om}$；本译本从页图。}\par\smallskip
\medskip\noindent\textbf{原书 p.582}\par
\noindent{\small\textbf{［概念］}“在平移意义下是层”指：等价 $\Phi_{\cE}$ 的核对象 $\cE$ 至多相差一次平移 $[n]$ 就是一个层。命题 5.2.2 把这件事归结为几何：$\cJ(\cE)$ 支撑在阿贝尔子簇 $X$ 上且等于 $i_*(L)$（定理 5.1.7），于是 $K=\Rb q'_*L$ 的上同调只落在长度不超过 $d=\dim\Ker q'$ 的区间内；再用 $K\cong\cE^\vee|_{(0,0)}\otimes\cE$ 逐个比较 $\cE$ 各上同调层的指标，就逼出只有一个非零层。原书特别指出，这一断言对 K3 曲面不成立。【来源：原书 p.581–582（断言、证明与“对 K3 曲面不成立”的说明）】}\par\smallskip
\medskip\noindent\textbf{原书 p.583}\par
\noindent{\small\textbf{［校勘］}页图此处印作 “Proposition 5.2.4”，而紧随其后的推论也印作 “Corollary 5.2.4”，编号重复。按 5.2.1、5.2.2 的顺序，本命题应为 5.2.3（本书引言亦引作“命题 5.2.3”）。为不改动原书面貌，此处照录原书编号，请读者注意。}\par\smallskip
\noindent{\small\textbf{［概念］}$\varphi_L$ 称为由线丛 $L$ 给出的同态：$T_d^*L\otimes L^{-1}$ 次数为零，故对应 $\hat D=\Pic^0(D)$ 的一个点；$L$ 丰富时 $\varphi_L$ 是同源，一般只是同态。嵌入 $\mathrm{NS}(D)\hookrightarrow\Hom(D,\hat D)$ 把除子类送到 $\varphi_L$，其像由自伴条件 $\hat\varphi=\varphi$ 刻画；下一段正是用 $\mathrm{NS}(D)\otimes\Q$ 来参数化 $D\times\hat D$ 上的对应。【来源：原书 p.583 正文；Mumford, \emph{Abelian varieties}, Oxford Univ. Press, 1970（本书文献 [32]）】}\par\smallskip
\medskip\noindent\textbf{原书 p.584}\par
\noindent{\small\textbf{［概念］}半齐次性把齐次性放松为“平移后至多差一个次数零的线丛”：齐次丛要求 $T^*_d\cE\cong\cE$，半齐次丛只要求 $T^*_d\cE\cong\cE\otimes L_d$ 且 $L_d\in\Pic^0(D)$。于是半齐次丛带有两个不变量：斜率 $\mu(\cE)$，以及 $\Pic^0$ 方向上的“平移缺陷”，后者恰好由对应 $\Phi_\mu\subset D\times\hat D$ 记录（命题 5.3.6）。命题 5.3.3 则说明齐次丛正是幂单丛被 $\Pic^0(D)$ 中线丛扭折后的直和。【来源：原书 定义 5.3.4、命题 5.3.3–5.3.6；[30]（Mukai, J. Math. Kyoto Univ. 18 (1978) 239–272）】}\par\smallskip
\medskip\noindent\textbf{原书 p.585}\par
\noindent{\small\textbf{［校勘］}本命题的陈述用 $\cF$ 表示被滤链的丛、用 $\cE_i$ 表示其子商，与定义 5.3.2 用 $\cF_i$ 表示滤链各项的记号不一致；这是原书自身的记号，照录。}\par\smallskip
\medskip\noindent\textbf{原书 p.586}\par
\noindent{\small\textbf{［校勘］}“另一方面，已知”一句原文只写 “it is known”，未给证明或出处，属跳步；按命题 5.3.6 的 (3) 与 $\Ker q_1$ 的构造可以补齐（译者核对）。}\par\smallskip
\noindent{\small\textbf{［校勘］}原书此处用独立的 “Construction 5.3.10” 环境；本书陈述环境表中没有 construction，故按体例“原文不用标准环境者作普通段落”处理，以粗体引出。另：$g$ 的矩阵项 $-\hat y^{-1}$ 指对偶同源 $\hat y:\hat B\to\hat A$ 之逆，原书未加说明，按上下文补出。}\par\smallskip
\noindent{\small\textbf{［概念］}$\Sigma^0(\cE)$ 是 $\cE$ 在 $\hat D=\Pic^0(D)$ 方向上的稳定化子：张量次数零线丛后仍同构于自身。引理 5.3.9 把它与对应 $\Phi_\mu$ 的竖直核联系起来：$\Sigma^0(\cE)=q_2(\Ker q_1)$，而由命题 5.3.6 的 (3)，$\deg q_1=r(\cE)^2$、$\deg q_2=\chi(\cE)^2$，两边计数便给出 $\operatorname{ord}\Sigma^0(\cE)=r^2$。“稳定化子的阶等于秩的平方”是后文判定两个纤维丛正交的关键。【来源：原书 式 (54)、引理 5.3.9、命题 5.3.6；[30] §7】}\par\smallskip
\medskip\noindent\textbf{原书 p.588}\par
\noindent{\small\textbf{［概念］}这条正合列把 $\Auteq\Db(A)$ 分成两层：离散部分 $\Z$ 来自平移 $[i]$，几何部分 $(A\times\hat A)_k$ 来自 $T_{a*}(\,\cdot\,)\otimes\cP_\alpha$；商 $U(A\times\hat A)$ 是 $A\times\hat A$ 的等距自同构群。序列不典范分裂，其障碍即末段所说的 2-上闭链（2-cocycle），显式公式见 [37]。于是“计算自等价群”被拆成两件可分别处理的事：算 $\Z\oplus(A\times\hat A)_k$，以及算中心扩张的类。【来源：原书 p.588 与 [37]（Orlov, Izv. Math. 66 (2002) 569–594）】}\par\smallskip
\medskip\noindent\textbf{原书 p.589}\par
\noindent{\small\textbf{［校勘］}末段中的 $\Ker(\varphi)$ 是 $\Ker(\varphi_L)$ 的简写（前几句均写 $\varphi_L$）；原书如此，照录。}\par\smallskip

\section*{参考文献（原书 p.591）}
\addcontentsline{toc}{section}{参考文献（原书 p.591）}

\medskip\noindent\textbf{原书 p.591}\par
\noindent{\small\textbf{［校勘］}原书 [6] 条目在文字层与页图上均印作 “Representablefunctors”，缺一词间空格；本译本按页图照录，未改。}\par\smallskip
\noindent{\small\textbf{［校勘］}原书 [44] 的 SGA 4 1/2 的分式在文字层被抽坏（(SGA 4 1 2 )），此处按页图排作 SGA 4½。}\par\smallskip
